\documentclass[11pt]{article}

% --- arXiv-ready preamble (no exotic packages; compiles with pdflatex) ---
\usepackage[utf8]{inputenc}
\usepackage[T1]{fontenc}
\usepackage[margin=1in]{geometry}
\usepackage{amsmath,amssymb,amsthm}
\usepackage{booktabs}
\usepackage{array}
\usepackage{natbib}
\bibliographystyle{abbrvnat}
\usepackage{hyperref}
\hypersetup{colorlinks=true,linkcolor=blue,citecolor=blue,urlcolor=blue}
\usepackage{microtype}
\usepackage{caption}

\title{No Pain, Same Gain: Interference-Gated Aggregation for
       Constant-Memory Recurrence}

\author{F\'abio de Oliveira Peres\\
        \small Independent Researcher, Rio Verde, Goi\'as, Brazil}

\date{Preprint, June 2026}

\begin{document}
\maketitle

\begin{abstract}
The Transformer established that fixed-state recurrence is dispensable for
sequence modeling---\emph{attention is all you need}---but attention pays for
that generality with a key--value cache that grows with context
(memory $\propto T$, $O(T^2)$ per-step cost). Growing-memory RNNs
\citep{behrouz2026memorycaching} reclaim recurrence by caching one associative
memory per segment so capacity grows with length, yet a faithful implementation
collapses on multi-query associative recall (MQAR): 25.2\% at 16 key--value pairs
versus 97.5\% for a \emph{fixed} linear-attention memory of the same width. We
trace the collapse to \emph{aggregation}, not routing: normalizing each cached
segment's readout before mixing destroys the global-denominator suppression that
makes linear attention work, so any non-saturated gate mixes full-scale garbage
at signal magnitude. We propose \textbf{Interference-Gated Aggregation (IGA)},
which gates \emph{unnormalized} per-segment numerators and denominators inside a
single global ratio; with gates initialized to one, IGA is numerically identical
to the fixed memory (verified to $5\times10^{-7}$), so the fixed memory is a
\emph{floor} and collapse is impossible by construction. We pair it with a
\textbf{tensor-sketch polynomial router} that routes by \emph{learned identity}
without reading token IDs---a differentiable relaxation of a count-min
sketch---giving recall that is flat in evaluation length even when queries never
repeat keys. We then make the recurrent state \textbf{constant in sequence
length}: a rank-$\le S$ factorization removes the $d^2$ term exactly, and capping
live segments with a low-rank compressed background bounds growth---$0.39$~MB at
$d=64$ versus a KV cache that grows $\propto T$, below it from
$T\!\approx\!1\mathrm{k}$ and $15\times$ smaller projected to a 1T-parameter model,
at no recall loss on the headline task. On the literature's own benchmark---the
official Zoology MQAR harness, hardest standard setting $(512,64)$---IGA reaches
$99.2\pm0.1$ versus Mamba $23.2\pm14.0$ and Based $79.5\pm1.4$, matching full
softmax attention ($100.0$). On a real character-level language model
(TinyStories, 2 layers, $d=128$, 3 seeds, two contexts) IGA matches a fixed-memory
baseline of the same width in bits-per-character within seed noise---\emph{constant
memory at no perplexity cost}---while a cost analysis shows the one-time training
premium is amortized by per-query inference savings that grow with context, since
IGA is $O(T)$ per step with a constant state where softmax attention is $O(T^2)$
per step with a $\propto T$ cache. \textbf{No pain, same gain}: the cost of
growing memory is removed, the quality is retained.
\end{abstract}

\section{Introduction}

``Attention is all you need'' \citep{vaswani2017attention} retired fixed-state
recurrence for sequence modeling, but the price is written into the mechanism: a
key--value cache that holds every past token. The cache grows linearly with
context and the per-step attention product is quadratic, so memory and compute
both scale with the sequence. This is the \emph{pain} of the attention
paradigm---bearable at short context, dominant at long context, and eventually
the binding constraint.

Growing-memory RNNs \citep{behrouz2026memorycaching} offer the opposite bet: keep a
recurrent state, but let it \emph{grow} by caching one associative memory per
segment of $S$ tokens and aggregating the cached readouts with a learned gate.
Capacity then scales with length while per-token cost stays sub-quadratic. The
promise is attention's recall with recurrence's cost structure. We find that a
faithful single-layer implementation of this promise does not merely fail to beat
the fixed memory it generalizes---it collapses far below it (Section~4, Table~1).

This paper makes a single, sharp proposal: a growing-memory model that carries
\textbf{no pain, same gain}---no growing-memory cost, and the same quality as the
strong baselines. We deliver it in three steps. First, we fix the aggregation: the
collapse is caused by per-segment readout normalization, and gating the
\emph{unnormalized} per-segment state inside one global ratio (IGA) removes it,
with the fixed memory as a provable initialization floor. Second, we make the
remaining bottleneck---segment \emph{routing}---solvable without token IDs, via a
tensor-sketch polynomial router that routes by learned identity. Third, we make
the state \emph{constant in sequence length}: a rank-$\le S$ factorization removes
the quadratic-in-width term exactly, and capping live segments with a low-rank
compressed background caps growth, giving a constant sub-KV state at no recall loss
on the headline task. The tagline names exactly what the measurements support:
\emph{no pain}---the growing memory and the $O(T^2)$ step are gone, leaving a
constant state and $O(T)$ per-step compute; \emph{same gain}---recall matches full
attention on the literature's hardest standard task ($99.2$ vs $100.0$) and
bits-per-character matches the fixed-memory baseline on real text (within seed
noise). The title is a compression of the result, not a slogan.

\paragraph{Contributions.}
\begin{enumerate}
\item \textbf{Diagnosis and aggregation fix (\S3.1).} The growing-memory collapse
  is an \emph{aggregation} failure: per-segment normalization erases the
  match-mass difference between segments. IGA gates unnormalized
  numerator/denominator pairs inside one global ratio; at gate$\approx 1$ (init)
  the model is the fixed memory to $5\times10^{-7}$, so learning can only
  \emph{remove} cross-segment interference---the source of the capacity that makes
  $K$ small memories beat one large memory.

\item \textbf{Learned identity routing (\S3.2).} Summed-similarity routing is
  kernel-family limited: it tracks the fixed memory's saturation across score
  families and even under auxiliary supervision, whereas identity-preserving
  summaries route exactly but exploit verbatim key repetition. Our tensor-sketch
  polynomial router \citep{pham2013tensorsketch} is a differentiable, ID-free
  relaxation of a count-min sketch that preserves the flat-recall property, and
  remains flat on \textbf{MQAR-indirect}, a benchmark where queries never repeat
  keys---recovered label-free by distilling a quadratic teacher trained only on
  the recall loss.

\item \textbf{Constant-memory state (\S3.3).} A rank-$\le S$ factorization of the
  per-segment cache removes the $d^2$ term exactly ($<\!5\times10^{-7}$); capping
  live segments at $K_{\max}$ with a low-rank compressed background makes the
  recurrent state constant in $T$---$0.39$~MB at $d=64$, below the KV cache from
  $T\!\approx\!1\mathrm{k}$ and $15\times$ smaller projected to 1T parameters,
  with no recall loss on the $(512,64)$ task.

\item \textbf{Real-text validation and deployment cost (\S5--\S6).} On a real
  character-level language model (3 seeds, two contexts) IGA matches a
  fixed-memory baseline of the same width in bits-per-character within seed
  noise---the constant-memory state carries no perplexity penalty---while the
  associative-recall advantage is shown to be task-specific (it does not transfer
  to a planted-fact probe on natural text at this scale). A cost analysis shows
  the one-time training premium (the gate) is amortized by per-query inference
  savings that grow with context, since IGA is $O(T)$ per step with a constant
  state against softmax attention's $O(T^2)$ per step with a $\propto T$ cache.
\end{enumerate}

We position the recall claims at small scale (single layer, $d=64$, synthetic
recall) and compare against exact attention, Based \citep{arora2024based} and
Mamba \citep{gu2023mamba} under the official Zoology protocol
\citep{arora2023zoology} with 3 seeds throughout; the language-model experiment
(\S5) is the one real-text setting.

\section{Background and setup}

\paragraph{MQAR.} Input is \texttt{k1 v1 \ldots kN vN SEP kq1 vq1 \ldots kqQ vqQ};
the model must predict each queried value as the next token after its query key.
We use the Zoology-style formulation with $N$ up to 128 pairs, $Q=4$ queries, key
vocabulary 192, value vocabulary 16 (chance 6.25\%); training always at $N=48$,
4000 steps, batch 256, AdamW $3\times10^{-4}$, identical seed before each model;
evaluation at $N\in\{48,64,96,128\}$ with no retraining. In \S4 we additionally use
the \emph{official} Zoology generator and training protocol (vocabulary 8192,
power-law query gaps, learning-rate sweep over $\mathrm{logspace}(-4,-2,4)$,
weight decay 0.1, cosine schedule, early stop at 99\%).

\paragraph{Fixed memory (FIXA).} Causal linear attention
\citep{katharopoulos2020linear} with feature map
$\varphi(x)=\mathrm{elu}(x)+1$:
\[
y_t = \frac{\sum_{j\le t}\varphi(q_t)\!\cdot\!\varphi(k_j)\,v_j}
            {\sum_{j\le t}\varphi(q_t)\!\cdot\!\varphi(k_j)}.
\]
A short depthwise causal convolution ($k=3$) before the projections lets one
layer bind key--value pairs, as in Based/Mamba. Capacity is bounded by $d^2$.

\paragraph{Growing memory (GRM).} Per segment $i$:
$M^{(i)}=\sum_j v_j\varphi(k_j)^\top$,
$z^{(i)}=\sum_j\varphi(k_j)$. At query $t$, each cached segment contributes a
\emph{normalized} readout
$r_i = M^{(i)}q_t / (z^{(i)}\!\cdot\!q_t)$, mixed by a softmax over match masses.
This is the mechanism that collapses.

\section{Method: Interference-Gated Aggregation}

\subsection{Aggregation that cannot collapse}

The fixed memory is \emph{exactly} the gate-$\equiv 1$ aggregation of
unnormalized per-segment caches:
\[
y_{\mathrm{fix}} = \frac{\sum_i M^{(i)}q}{\sum_i z^{(i)}\!\cdot\!q}.
\]
IGA inserts the gate inside this ratio:
\[
y_t = \frac{N_{\mathrm{cur}} + \sum_{i<s} g_i\,M^{(i)}q_t}
            {m_{\mathrm{cur}} + \sum_{i<s} g_i\,z^{(i)}\!\cdot\!q_t},
\qquad
g_i = \sigma\!\big(a(\log s_i - \mathrm{ref}) + b\big),
\]
where $s_i$ is a routing score and $\mathrm{ref}$ is the mean (similarity scores)
or max (count scores) of $\log s_i$. With $b=3$, $g\approx 0.95$ everywhere at
init $\Rightarrow$ the model \emph{starts} as the fixed memory; the unit test
\texttt{test\_ungated\_equals\_fixed} verifies the identity to $5\times10^{-7}$.
The gate can only \emph{remove} interference---which is precisely where the
capacity of $K$ small memories over one large memory comes from.

\paragraph{Why GRM collapses.} Instrumenting the trained GRM at $n=16$: recall
$25.2\%$, routing accuracy of its softmax gate $90.7\%$, gate mass on the current
segment $\approx 0$. The killer is that normalized readouts from wrong segments
are garbage \emph{at the same magnitude as the signal}; a $0.7/0.3$ softmax mixes
$30\%$ full-scale garbage. No fix to the softmax can help; the normalization must
move \emph{outside} the mixture---which is what IGA does.

\subsection{Learned identity routing}

With aggregation fixed, IGA's residual gap to oracle is pure routing. The routing
problem is a \emph{kernel-family} problem: every summed-similarity score
(first/second moment, elementwise-exponential features, learned embeddings with
auxiliary supervision) saturates exactly like the fixed memory, whereas
\emph{identity-preserving} summaries---the token histogram or its count-min
sketch \citep{cormode2005countmin}---route exactly (100\%/93\%+) at zero
parameters and give recall flat in evaluation length, but they read token IDs and
MQAR queries repeat keys verbatim. Real language does not.

We need a score $s_i(q)\approx$ (occurrences of $q$'s identity in segment $i$),
cacheable as a per-segment sum, differentiable, and ID-free. For near-orthonormal
embeddings $\hat e$,
$(\hat e_q\!\cdot\!\hat e_k)^p\to\delta(q=k)$ as $p$ grows, and the degree-$p$
polynomial kernel is exactly linearizable:
$(x\!\cdot\!y)^p = \langle x^{\otimes p},y^{\otimes p}\rangle$, compressible by
tensor sketch \citep{pham2013tensorsketch} with
$\mathbb{E}[\mathrm{ts}(x)\!\cdot\!\mathrm{ts}(y)]=(x\!\cdot\!y)^p$ and
$O(1/m)$ variance for unit vectors---\emph{polynomial} cost in sharpness, where
Performer-style estimates of a comparably sharp $\exp(\tau q\!\cdot\!k)$ require
variance exponential in $\tau$ \citep{choromanski2021performer}. The router score
is
\[
s_i = \mathrm{softplus}(w_1)\big(\hat e_q\!\cdot\!\textstyle\sum_j\hat e_{k_j}\big)
      + \mathrm{softplus}(w_p)\,\mathrm{med}_{J}\!\big(\mathrm{ts}(\hat e_q)\!\cdot\!\textstyle\sum_j\mathrm{ts}(\hat e_{k_j})\big),
\]
with three design elements: (i) the \textbf{median of $J=3$ independent
sketches}---the continuous analogue of count-min's min---lifts routing from
$96.6\%$ to $99.8\%$ in a no-training microbenchmark on the exact MQAR segment
statistics ($p=4$, $m=256$, $r=64$), where degree-1 alone yields only
$39.9$--$60.8\%$ (the same saturation that limits the fixed memory); (ii) a
\textbf{degree-1 term for cold start}, since
$\partial(q\!\cdot\!k)^p/\partial\mathrm{emb}$ vanishes at orthogonality for
$p\ge 2$, giving word2vec-style alignment gradients when embeddings must be
\emph{learned}, with learnable mixture weights that let training anneal it away;
(iii) \textbf{count-vs-similarity gate links}---count scores gate on
$\log(\mathrm{clamp}(s))$ with a max reference (hard suppression from init), while
learned-poly scores gate on $\log(\mathrm{softplus}(s))$ and supervise raw scores
as logits, because learned scores can be negative at cold start and
$\log(\mathrm{clamp})$ creates a dead gradient zone that destroys readout training
(this distinction cost one full failed run, documented as a negative result).
Summary cost: $r + J\!\cdot\!m = 832$ floats/segment (count-min: 256; histogram:
$V$).

\paragraph{MQAR-indirect: identity must be learned.} We introduce a benchmark in
which the query block uses a \emph{disjoint} query vocabulary, paired to keys by a
bijection fixed across the dataset---sequence structure, positions and protocol
unchanged. By construction, ID-counting routers receive zero signal. Two training
regimes: (a) auxiliary routing cross-entropy with bijection labels---\emph{privileged}
supervision, an upper bound for the kernel family; (b) \emph{distillation}---a
minimal quadratic softmax-attention teacher ($d=64$, 1 layer, own embeddings)
trained only on the recall loss, whose per-segment attention mass is distilled
into the router by KL; the teacher is discarded at inference (quadratic cost only
during training, mirroring SMT \citep{kumar2026smt}). The teacher solves the task
at $100\%$, confirming softmax attention does not suffer the perverse equilibrium.

\subsection{Constant-memory state}

A growing memory that beats the fixed memory on recall but grows without bound
does not yet answer the paradigm question---its state would eventually exceed the
KV cache it is meant to replace. We make the recurrent state \textbf{constant in
$T$} in two steps, both presented as design rather than patch.

\paragraph{Factored state.} The per-segment cache
$M^{(i)}=\sum_{j\in\mathrm{seg}_i}v_j\varphi(k_j)^\top$ is a sum of $S$ rank-1
updates, so $\mathrm{rank}(M^{(i)})\le S=16$. We factor it exactly as
$M^{(i)}={U^{(i)}}^\top W^{(i)}$ with $U^{(i)},W^{(i)}\in\mathbb{R}^{S\times d}$
(rank $R=S$), storing $2dS$ floats per segment instead of $d^2$ and computing the
readout $M^{(i)}q = {U^{(i)}}^\top(W^{(i)}q)$ at $2dS$ cost instead of $d^2$. The
factorization is exact for a rank-$\le S$ matrix, so the model output is
unchanged up to floating-point error---\texttt{test\_factored\_equals\_dense}
verifies the identity (threshold $10^{-4}$; measured gap $<\!5\times10^{-7}$)---and
at the official Zoology task $(512,64)$ the factored model matches the dense one
within seed noise (dense $99.02\pm0.17$, factored $99.18\pm0.08$, 3 seeds).
Factoring removes the quadratic-in-width term ($d^2\to 2dS$: $4096\to 2048$
elements at $d=64$), but the state still grows with the number of segments,
$O(T/S)$.

\paragraph{Capped memory.} We bound the number of \emph{live} cached segments at
$K_{\max}$ and compress the overflow into a single low-rank background of rank
$r_{\mathrm{bg}}$. The recurrent state becomes
\[
\mathrm{state} = K_{\max}\!\cdot\!\mathrm{per\_seg}
               + (2d\,r_{\mathrm{bg}} + d + m) + K_{\max},
\]
independent of $T$, where $\mathrm{per\_seg}=2dS+d+m$. Compression is a
deterministic FIFO eviction with normalization of the compressed readout; live
segments and the compressed background are mixed inside IGA's single global ratio,
so the interference-gated aggregation of \S3.1 is preserved. At the official task
$(512,64)$ the context is exactly $K_{\max}=32$ segments, so no eviction occurs and
the capped model is a no-op: recall is identical to the uncapped factored
model---$99.23/99.30/99.02$ (mean $99.18$) across three seeds. The honest price of
constant memory is that segments older than $K_{\max}$ are compressed (lossy):
recall of facts within the last $K_{\max}$ segments tracks the unlimited baseline
even with hundreds of evictions, while older facts are not individually
recoverable---the explicit recall/memory tradeoff.

\section{Experiments---associative recall}

All numbers are mean $\pm$ std over 3 seeds unless noted. Baselines: exact
attention (2-layer MHA with learned positions), Based (official hybrid: BaseConv +
Taylor-exp linear attention), Mamba (S6; pure-PyTorch reimplementation with exact
parallel scan, unit-tested against the sequential reference). Backbone,
initialization (GPT-2 scheme) and hyperparameters follow the Zoology Figure-2
configuration. A negative finding, verified against the \emph{official}
implementation: the 2-layer attention-only Figure-2 configuration does not form the
induction circuit under our infinite-stream protocol (fresh samples every step
remove the memorization bootstrap); it plateaus at $\sim$$13\%$ across learning
rates, and the official \texttt{zoology.model.LanguageModel} trained under the same
loop gives $12.3\%$---identical to our replica---while under the official
fixed-data regime our replica \emph{does} reproduce the published induction-head
transition. We therefore report \textbf{attention-conv}---a single-layer
exact-attention model with the same short-convolution stem every other model
has---as the working upper anchor.

\subsection{Verbatim MQAR (our protocol)}

Train at $n=48$ ($n_k=192$), evaluate without retraining; chance 6.25\%.

\begin{table}[h]\centering
\begin{tabular}{lcccc}
\toprule
model & $n{=}48$ & $n{=}64$ & $n{=}96$ & $n{=}128$ \\
\midrule
attention (Zoology fig2) & 12.9\,$\pm$\,0.0 & 11.1\,$\pm$\,0.5 & 9.7\,$\pm$\,0.5 & 9.0\,$\pm$\,0.4 \\
based                   & 8.3\,$\pm$\,0.1  & 7.9\,$\pm$\,0.2  & 7.3\,$\pm$\,0.2 & 7.3\,$\pm$\,0.1 \\
mamba                   & 13.3\,$\pm$\,0.0 & 12.3\,$\pm$\,0.1 & 10.6\,$\pm$\,0.2 & 9.6\,$\pm$\,0.4 \\
FIXA(64)                & 77.9\,$\pm$\,1.0 & 71.6\,$\pm$\,0.8 & 59.9\,$\pm$\,0.4 & 50.9\,$\pm$\,0.4 \\
FIXA(128)               & 92.1\,$\pm$\,0.5 & 88.8\,$\pm$\,0.6 & 81.8\,$\pm$\,1.0 & 74.5\,$\pm$\,1.1 \\
IGA-hist (IDs)          & 98.3\,$\pm$\,0.1 & 98.3\,$\pm$\,0.1 & 98.4\,$\pm$\,0.2 & \textbf{98.5\,$\pm$\,0.1} \\
IGA-sketch (IDs)        & 98.1\,$\pm$\,0.1 & 98.1\,$\pm$\,0.1 & 98.1\,$\pm$\,0.2 & 98.0\,$\pm$\,0.2 \\
\textbf{IGA-poly (no IDs)} & 97.7\,$\pm$\,0.1 & 97.6\,$\pm$\,0.2 & 97.5\,$\pm$\,0.2 & \textbf{97.1\,$\pm$\,0.2} \\
\bottomrule
\end{tabular}
\end{table}

All IGA variants are flat in evaluation length; the ID-free polynomial router
pays $0.8$--$1.4$ points of sketch noise relative to the exact histogram. The
2-layer literature baselines fail in this regime (48 pairs at $d=64$, far beyond
their capacity-favorable setting).

\subsection{MQAR-indirect}

Same protocol; the query never repeats the key.

\begin{table}[h]\centering
\begin{tabular}{lcccc}
\toprule
model & $n{=}48$ & $n{=}64$ & $n{=}96$ & $n{=}128$ \\
\midrule
attention (Zoology fig2) & 13.1\,$\pm$\,0.1 & 11.5\,$\pm$\,0.4 & 9.9\,$\pm$\,0.2 & 9.2\,$\pm$\,0.2 \\
based                     & 7.4\,$\pm$\,0.3  & 7.3\,$\pm$\,0.1  & 6.8\,$\pm$\,0.2 & 6.8\,$\pm$\,0.1 \\
mamba                    & 13.3\,$\pm$\,0.0 & 12.2\,$\pm$\,0.1 & 10.6\,$\pm$\,0.3 & 9.6\,$\pm$\,0.3 \\
FIXA(64)                 & 57.6\,$\pm$\,20.8 & 51.4\,$\pm$\,19.4 & 42.5\,$\pm$\,16.2 & 36.8\,$\pm$\,13.9 \\
IGA-sketch (control)     & 57.2\,$\pm$\,20.6 & 50.9\,$\pm$\,19.1 & 42.1\,$\pm$\,15.9 & 36.4\,$\pm$\,13.6 \\
IGA-emb+aux (exp family) & 68.6\,$\pm$\,17.8 & 62.2\,$\pm$\,17.9 & 52.1\,$\pm$\,16.0 & 45.1\,$\pm$\,13.8 \\
\textbf{IGA-poly+aux (privileged)}        & 96.1\,$\pm$\,0.7 & 95.8\,$\pm$\,0.7 & 95.4\,$\pm$\,0.7 & \textbf{94.8\,$\pm$\,0.9} \\
\textbf{IGA-poly distilled (label-free)}   & 93.5\,$\pm$\,1.0 & 92.6\,$\pm$\,1.4 & 90.5\,$\pm$\,1.8 & \textbf{88.2\,$\pm$\,2.2} \\
\bottomrule
\end{tabular}
\end{table}

The ID-counting control tracks FIXA seed-by-seed (routing $6.0\%=$ chance),
confirming the benchmark removes the verbatim shortcut by construction. The
polynomial-kernel router preserves the flat-recall property with \emph{fully
learned} identity: $+58$ points over the fixed memory at $n=128$, and the
label-free distilled variant retains $93$--$96\%$ of the privileged upper bound.

\subsection{Official Zoology protocol (vocabulary 8192)}

Official generator and training loop; 3 seeds, mean $\pm$ std.

\begin{table}[h]\centering
\begin{tabular}{lcc}
\toprule
model & $(128,8)$ easy & $(512,64)$ hard \\
\midrule
attention (full softmax)          & 99.4\,$\pm$\,0.3 & 100.0\,$\pm$\,0.0 \\
attention-conv (1-layer anchor)   & 99.3\,$\pm$\,0.3 & 99.4\,$\pm$\,0.0 \\
based                             & 99.2\,$\pm$\,0.2 & 79.5\,$\pm$\,1.4 \\
mamba                             & 99.1\,$\pm$\,0.0 & 23.2\,$\pm$\,14.0 \\
\textbf{IGA-sketch (ours)}        & 99.1\,$\pm$\,0.0 & \textbf{99.1\,$\pm$\,0.0} \\
\textbf{IGA-poly (ours, learned identity)} & 99.1\,$\pm$\,0.0 & \textbf{99.2\,$\pm$\,0.1} \\
\bottomrule
\end{tabular}
\end{table}

The discriminating regime is $(512,64)$: 64 key--value pairs over a 512-token
context ($32$ segments of $S=16$---$2\times$ more segments than any setting
above). All six models solve the easy task. On the hard task the efficient
baselines separate exactly as the Based paper reports: \textbf{Mamba collapses
($23.2\pm14.0$, one seed at $7\%$---the fixed-state saturation) and Based degrades
($79.5\pm1.4$).} \textbf{Both IGA models hold at $\sim\!99\%$}, matching full
softmax attention ($100\%$) and the convolutional attention anchor ($99.4$), at
sub-quadratic per-token cost. This is the citable result that closes the
comparison: on the literature's own benchmark, in its hardest standard setting,
the growing memory keeps near-perfect recall where the fixed-state models fail.

\subsection{Inference cost}

Batch 1, $d=64$, fp32, reference Python loops on both sides (scaling comparison,
not engineering constants).

\begin{table}[h]\centering
\begin{tabular}{lcccl}
\toprule
model & $T{=}4096$ & $T{=}16384$ & $T{=}65536$ & state @ $T{=}65536$ \\
 & tok/s & tok/s & tok/s & \\
\midrule
attention (SDPA, KV cache)      & 3577 & 997 & 253 & 32.0~MB ($\propto T$) \\
IGA-poly recurrent, soft gate   & 3337 & 3721 & 3167 & \textbf{0.39~MB} (const) \\
IGA-poly recurrent, hard top-1 & 4885 & 4725 & \textbf{4455} & \textbf{0.39~MB} (const) \\
\bottomrule
\end{tabular}
\end{table}

IGA throughput is $\sim$flat in context length; incremental exact attention
degrades linearly, $17.6\times$ slower at $64$k. The state column reports the
constant-memory state of \S3.3---$0.39$~MB constant---against a KV cache that
grows $\propto T$ ($32$~MB at $64$k). The throughput crossover and the constant
sub-KV state are the claims.

\subsection{Constant-memory scaling and the long-context regime}

With the factored + capped state, the recurrent state is constant in $T$. Measured
with \texttt{state\_bytes} ($d=64$, $K_{\max}=32$, $r_{\mathrm{bg}}=16$, fp32) at
$0.39$~MB against the KV cache $2Td$:

\begin{table}[h]\centering
\begin{tabular}{lccl}
\toprule
$T$ & KV cache & IGA (constant) & ratio \\
\midrule
512   & 0.25~MB  & 0.39~MB & 1.53$\times$ \\
1024  & 0.50~MB  & 0.39~MB & 0.76$\times$ \\
2048  & 1.00~MB  & 0.39~MB & 0.38$\times$ \\
4096  & 2.00~MB  & 0.39~MB & 0.19$\times$ \\
65536 & 32.0~MB  & 0.39~MB & 0.012$\times$ \\
\bottomrule
\end{tabular}
\end{table}

The state equals the cache at the crossover $T\approx 2K_{\max}S$ (where the cache
fills exactly $K_{\max}$ segments); below it the constant state is larger---it is
sized for capacity, not the current context---and above it the advantage grows
without bound. Projected to a 1T-parameter regime ($d=8192$, $K_{\max}=4096$,
$r_{\mathrm{bg}}=4$), the constant state is $209.5$~GB $=0.065\times$ the $3.1$~TB
KV cache---\textbf{$15\times$ smaller, and constant}---where the dense IGA state
would be $800$~TB ($256\times$ larger) and even the factored-but-uncapped state is
$3.2$~TB (parity, still growing). The constant sub-KV state is the memory result.

\section{Real-text language modeling: constant memory at no perplexity cost}

All \S4 results use synthetic associative recall. We close the remaining
question---does the constant-memory IGA carry a perplexity penalty on \emph{real}
text?---by training a 2-layer, $d=128$ IGA language model (segment size 64;
tensor-sketch polynomial router reading raw tokens; factored state of \S3.3)
against the \emph{identical} architecture with the gate disabled
(\texttt{gated=False}, the FIXA fixed-memory baseline of the same width and
parameter count). Data is character-level TinyStories, 88~MB, downloaded---no
synthetic data is generated for any language-model experiment. Optimizer AdamW
$3\times10^{-4}$, batch 16, 4000 steps, 3 seeds, contexts $\{1024, 2048\}$. We
report bits-per-character (bpc) on a 5\% held out split and a planted-fact recall
probe---the MQAR analogue on natural text---at three distances (ctx/4, ctx/2,
ctx$-$64).

\begin{table}[h]\centering
\begin{tabular}{lccccc}
\toprule
ctx & IGA bpc & FIXA bpc & $\Delta$ bpc (IGA$-$FIXA) & recall@ctx$-$64 IGA & recall@ctx$-$64 FIXA \\
\midrule
1024 & 2.1607\,$\pm$\,0.0060 & 2.1621\,$\pm$\,0.0084 & $-$0.0015 & 1.67 & 1.67 \\
2048 & 2.1138\,$\pm$\,0.0114 & 2.1143\,$\pm$\,0.0127 & $-$0.0005 & 2.29 & 2.08 \\
\bottomrule
\end{tabular}
\end{table}

Two findings. (i) \textbf{No perplexity cost.} The bpc difference is $-0.0015$
(ctx 1024) and $-0.0005$ (ctx 2048), both an order of magnitude below the per-seed
std ($0.006$--$0.013$); IGA is never consistently worse across the six (model,
seed) runs. The constant-memory state therefore does not degrade language
modeling relative to a fixed-memory model of the same width and parameter count.
(ii) \textbf{The recall advantage is task-specific.} The planted-fact probe is
near chance for both models at both contexts ($1.7$--$2.3\%$, per-seed spread
$0$--$4.4\%$---pure noise, no distance dependence); the flat-recall advantage of
\S4 does not transfer to retrieval of an arbitrary planted fact in natural text at
$d=128$. The structured MQAR benchmark remains the right instrument for the
associative-recall claim; natural-text retrieval at this scale does not separate
the architectures, and scaling $d$/ctx to surface a separation is left to future
work (GPU-budget-bound).

\paragraph{Compute.} IGA costs $5.2\times$ (ctx 1024) and $4.9\times$ (ctx 2048)
more wall-clock per step than FIXA in our reference Python-loop implementation---the
gate's per-segment normalization, dominated by loop overhead rather than
asymptotic FLOPs. The tradeoff is memory-vs-compute (constant $O(1)$ state against
higher per-step compute), not memory-vs-quality: quality is matched within noise.

We note honestly that a preliminary 300-step run in the parent project reported a
$\sim$$0.016$ bpc gain at ctx $\ge 2048$; it does \textbf{not} reproduce under the
3-seed, 4000-step protocol here ($\Delta=-0.0005$, within noise). The
language-model result is a no-regression result, not a quality gain.

\section{Deployment economics: where constant memory pays}

This section synthesizes \S4.4 (throughput), \S4.5 (memory) and \S5 (LM quality)
into the deployment-cost picture---the regime in which a growing-memory model is
actually meant to replace a KV-cache transformer.

\paragraph{Training (one-time).} Per step, softmax attention is $O(T^2 d)$---the
QK$^\top$ and AV products over the full $T\!\times\!T$ matrix---while IGA's
linear-attention recurrence is $O(Td^2)$ (the factored state update is $O(TdS)$,
smaller). The asymptotic FLOP scaling therefore favors IGA as context grows, with
ratio $T/d$: $\sim$$8\times$ fewer attention-FLOPs at ctx 1024 ($d=128$),
$\sim$$512\times$ at $T=64$k ($d=64$). We did \emph{not} train a transformer
language-model baseline, so we state the training-cost comparison as complexity
analysis, not as a measured number; the one comparison we \emph{do} measure
end-to-end is inference throughput against exact attention (\S4.4: $14$--$17.6\times$
at $64$k). The $5\times$ wall-clock premium of \S5, in contrast, \emph{is}
measured---against the \emph{fixed-memory linear baseline} (FIXA), the relevant
same-family control for the gate's overhead.

\paragraph{Inference (recurring, per generated token).} A KV-cache transformer
reads the full cache per token---$O(TdL)$ memory traffic---and holds a cache of
$2TdL$ that grows $\propto T$. IGA reads a constant $O(K_{\max}dS)$ per token and
holds a constant state. Measured at batch 1, $d=64$ (\S4.4): IGA holds
$\sim$$4.5$k tok/s flat from $T=4$k to $64$k, while incremental exact attention
degrades $\sim$$14\times$ to $253$ tok/s at $64$k. Per-token inference cost of IGA
is constant in $T$; the transformer's grows with $T$.

\paragraph{Amortization.} Training is paid once; inference is paid per query
across the model's serving lifetime. The one-time $5\times$ gate premium (against
the linear baseline) is amortized by the recurring inference savings once the
served context is long enough that the constant state undercuts the cache---the
crossover $T\approx 2K_{\max}S$ of \S4.5, $15\times$ smaller projected to 1T. Below
the crossover, where the cache is itself small and attention is fast, the
transformer is the cheaper option and IGA's advantage is only the recall capacity
of \S4; above it, IGA's per-token cost stops growing while the transformer's keeps
growing, so the total cost of serving long contexts eventually favors IGA by an
unbounded margin. Past the point where the KV cache no longer fits on device, IGA
is not the cheaper option but the only feasible one.

This is a cost-structure argument, not a measured deployment dollar figure:
absolute cost depends on hardware, batching and kernel fusion not modeled here
(reference Python loops). The structural claim---$O(1)$ memory and
$O(1)$-in-$T$ per-token compute against $O(T)$ and $O(T)$ for the cache, with
matched LM quality (\S5)---is what the constant-memory result buys. That is the
``no pain'': the growing-memory cost is removed. The ``same gain'' is the matched
recall (\S4.3) and matched perplexity (\S5).

\section{Related work}

\paragraph{Efficient-attention recall benchmarks.} Zoology
\citep{arora2023zoology} established MQAR as the discriminating task between
attention and efficient mixers; Based \citep{arora2024based} frames the
recall-throughput tradeoff. We adopt their harness and add MQAR-indirect, which
isolates \emph{identity learning} from \emph{identity exploitation}---to our
knowledge the first recall benchmark where queries never repeat keys by
construction.

\paragraph{Growing memory.} Memory Caching \citep{behrouz2026memorycaching}
proposes segment-cached memories with gated aggregation (GRM among its variants).
We keep its state layout ($M^{(i)}, z^{(i)}$) and cost, and replace its aggregation
(the collapse cause) and its routing (the saturation cause). Its state grows with
the number of cached segments ($O(T/S)$, quadratic in width via $d^2$); our
factored + capped state (\S3.3) is constant in $T$ and linear in width, removing
the growth that would otherwise exceed the KV cache it is meant to replace.

\paragraph{Polynomial sketches in attention.} PolySketchFormer
\citep{kacham2023polysketchformer} uses polynomial sketches to \emph{replace the
attention kernel itself}. Our use is different: the sketch summarizes
\emph{segment identity} for routing in a growing-memory RNN; attention/readout
keeps its own feature map. Count sketch \citep{charikar2002countsketch}, count-min
\citep{cormode2005countmin} and tensor sketch \citep{pham2013tensorsketch} supply
the estimators; the median-of-$J$ variant mirrors count-min's min under signed
estimates.

\paragraph{Linear attention and SSMs.} The fixed memory is
\citet{katharopoulos2020linear}; Based and Mamba are the strongest small-scale
recall baselines in this family. Performer \citep{choromanski2021performer}
approximates $\exp(q\!\cdot\!k)$ with random features; our variance argument
(\S3.2) explains why that family cannot provide $\delta$-sharp identity routing at
fixed feature budget.

\paragraph{Distillation across compute classes.} SMT \citep{kumar2026smt} trains
an RNN with a parallel teacher; our teacher-to-router distillation transfers
\emph{routing} (not state updates) and is, to our knowledge, the first label-free
recovery of identity routing in a growing-memory model.

\section{Limitations and honest notes}

\begin{itemize}
\item Single-layer $d=64$ synthetic recall is the headline setting; the language
  model of \S5 uses 2 layers, $d=128$, char-level TinyStories, 3 seeds, two
  contexts (downloaded corpus, no synthetic data). IGA matches the fixed-memory
  baseline in bpc within seed noise (no perplexity cost), but the planted-fact
  recall probe is near chance for both at this scale, and the $\sim$$0.016$ bpc
  gain reported by a preliminary 300-step run does \emph{not} reproduce under the
  3-seed 4000-step protocol; the associative-recall advantage is demonstrated on
  structured MQAR (\S4), not on natural-text retrieval at $d=128$. Scaling $d$ and
  context to surface a recall separation on real text is left to future work
  (GPU-budget-bound).
\item The indirect bijection is 1-to-1 and deterministic; real paraphrase is
  many-to-many and noisy.
\item The privileged-aux arm uses bijection labels; the distillation arm removes
  them but adds a quadratic teacher \emph{during training only}. No arm learns
  routing from the recall gradient alone---the routing-family analysis of \S3.2
  suggests this is impossible in this regime (perverse equilibrium); we state it
  as a conjecture, not a theorem.
\item Mamba is a pure-PyTorch reimplementation (exact scan, unit-tested); absolute
  wall-clock comparisons to fused-kernel implementations would be unfair and are
  not made---only task accuracy is compared.
\item The cost benchmark uses reference Python loops on both sides (no fused
  kernels); it measures \emph{scaling}, not engineering constants. The recurrent
  state is constant in $T$ (\S3.3) and below the KV cache from $T\approx 1$k
  onward; at the smallest contexts ($T\lesssim 2K_{\max}S$) the constant state is
  sized for capacity and exceeds the cache. Segments older than $K_{\max}$ are
  compressed (lossy), the explicit recall/memory tradeoff.
\item Seeds: 3 per cell. Std is sample std (Bessel). No hyperparameter search was
  performed for our models beyond the E0 microbenchmark ($p,m,J,r$); baselines
  received the official learning-rate sweep.
\end{itemize}

\section{Reproduction}

\begin{verbatim}
pip install -r requirements.txt
python3 tests/test_sanity.py          # 14 equivalence/property tests
                                       # (incl. factored==dense <5e-7, IGA==fixed at init)
python3 scripts/run_cost.py            # inference cost (~5 min)
python3 scripts/run_multiseed.py all   # our models, 3 seeds
python3 scripts/run_baselines_ours.py all   # literature baselines, our protocol
python3 scripts/run_zoology.py all     # official protocol (incl. (512,64))
python3 scripts/aggregate.py          # tables (mean +/- std)
python3 scripts/run_factored_check.py # factored vs dense recall at (512,64), 3 seeds
python3 scripts/run_fixb_check.py      # capped-memory recall == factored + constant-state table
python3 scripts/project_memory.py     # memory projection (small/medium/1T)
# Real-text language modeling (downloaded TinyStories; isolated folder):
python3 multicamada/run_lm_ml.py --steps 4000 --ctx 1024 --seeds 0,1,2 --tag lm_full
python3 multicamada/run_lm_ml.py --steps 4000 --ctx 2048 --seeds 0,1,2 --tag lm_full_ctx2048
\end{verbatim}

The language-model experiment uses the repository's existing multi-layer model
\texttt{IGALM} (\texttt{src/models/iga\_layer.py},
$n_{\mathrm{layers}}{=}2$, $d=128$, factored-capable) compared against itself with
\texttt{gated=False} (the FIXA baseline of matched width/parameters), isolated in
the \texttt{multicamada/} folder. The corpus is the downloaded TinyStories sample
(88~MB); no synthetic LM data is generated. All runs are resumable
(skip-if-exists per JSON in \texttt{results/raw/}).

\bibliography{references}

\end{document}